\Article{244}
If a form $f$ can represent the number $a$, a form $f'$ the number $a'$, and the form $F$ can be transformed into $ff'$, it is easy to see that the product $aa'$ can be represented by the form $F$. It follows at once that, when the determinants of these forms are negative, the form $F$ is positive if $f$, $f'$ are both positive or both negative; by contrast, $F$ is negative if one of the forms $f$, $f'$ is positive and the other negative. Let us focus especially on the case considered in the preceding article, where $F$ is composed from $f$, $f'$. Also, $f$, $f'$ and $F$ have the same determinant $D$. Suppose further that the representations of the numbers $a$, $a'$ by the forms $f$, $f'$ use coprime values of the indeterminates, and that the former belongs to the value $b$ of the expression $\surd D\ (\mathrm{mod.}\ a)$, the latter to the value $b'$ of the expression $\surd D\ (\mathrm{mod.}\ a')$, and set $bb-D=ac$, $b'b'-D=a'c'$. Then, by Article 168, the forms $(a,b,c)$, $(a',b',c')$ are properly equivalent to the forms $f$, $f'$, so $F$ is also composed from those two forms. But the form $(A,B,C)$ is composed from the same forms if, denoting the greatest common divisor of $a$, $a'$, $b+b'$ by $=\mu$, we set $A=\frac{aa'}{\mu\mu}$, $B\equiv b$ and $b'$ modulo $\frac{a}{\mu}$, $\frac{a'}{\mu}$. respectively,
\SourceLeaf{272}
$AC=BB-D$; so this form is properly equivalent to the form $F$. Now the form $Axx+2Bxy+Cyy$ represents the number $aa'$ on setting $x=\mu$, $y=0$, the greatest common divisor of these values being $\mu$. Thus $aa'$ can also be represented by the form $F$ so that the values of the indeterminates have greatest common divisor $\mu$ (Article~166). Whenever $\mu=1$, therefore, $aa'$ can be represented by the form $F$ by assigning coprime values to the indeterminates, and this representation belongs to the value $B$ of the expression $\surd D$ (modulo $aa'$) that is congruent to $b$, $b'$ modulo $a$, $a'$ respectively. The condition $\mu=1$ always holds when $a$, $a'$ are coprime; more generally, when the greatest common divisor of $a$, $a'$ is coprime to $b+b'$.

\Needspace{11\baselineskip}
\begin{center}
\textit{Composition of orders.}
\end{center}

\Article{245}
\textsc{Theorem.} \textit{If the form $f$ belongs to the same order as $g$, and likewise $f'$ to the same order as $g'$, the form $F$ composed from $f$, $f'$ has the same determinant and belongs to the same order as the form $G$ composed from $g$, $g'$.}

\textit{Proof.} Let the forms be $f,f',F=(a,b,c)$, $(a',b',c')$, $(A,B,C)$ respectively, and their determinants $=d,d',D$. Further, let the greatest common divisor of $a,2b,c$ be $=m$, and the greatest common divisor of $a,b,c$ be $=\mathfrak{m}$; let $m'$, $\mathfrak{m}'$ have similar meanings for the form $f'$, and $M$, $\mathfrak{M}$ for the form $F$. Then the order of the form $f$ is determined by the numbers $d,m,\mathfrak{m}$, so the same numbers also apply to the form $g$. Likewise, the numbers $d',m',\mathfrak{m}'$ are the same for the form $g'$ as for the form $f'$. Now, by Article~235, the numbers $D,M,\mathfrak{M}$ are determined by $d,d',m,m',\mathfrak{m},\mathfrak{m}'$: namely, $D$ is the greatest common divisor of $dm'm'$, $d'mm$; $M=mm'$; and $\mathfrak{M}=\mathfrak{m}\mathfrak{m}'$ (if both $m=\mathfrak{m}$, $m'=\mathfrak{m}'$) or $=2\mathfrak{m}\mathfrak{m}'$ (if $m=2\mathfrak{m}$ or $m'=2\mathfrak{m}'$). These properties of $D,M,\mathfrak{M}$ follow from the fact that $F$ is composed from $f,f'$, so it is easy to see that $D,M$ and $\mathfrak{M}$ also apply to the form $G$, and therefore $G$ belongs to the same order as $F$. \textit{Q.~E.~D.}

For this reason, we call the order containing the form $F$ composed from the orders containing the forms $f,f'$. Thus, \textit{for example,}\ two properly primitive orders always give an order of the same kind; a properly primitive and an improperly primitive order give an improperly primitive order. --- A similar meaning is intended when an order is said to be composed from several other orders.
